\documentclass[a4paper, 12pt]{article}
\usepackage[utf8]{inputenc}
\usepackage{amsmath, amssymb, amsfonts}
\usepackage{geometry}
\usepackage{microtype} 
\usepackage{pgfplots}
\usepackage[spanish, es-noshorthands]{babel} 
\usepackage{setspace}
\usepackage{tikz}
\usepackage{hyphenat}
\usepackage[hidelinks]{hyperref}
\usepackage{fancyhdr}

\usetikzlibrary{patterns}
\pgfplotsset{compat=1.17} 
\usetikzlibrary{decorations.markings}

% --- CONFIGURACIÓN DE MÁRGENES ---
\geometry{
 a4paper,
 left=2.5cm,
 right=2.5cm,
 top=2.5cm, 
 bottom=2.5cm,
 includehead
}
\setlength{\parindent}{0pt}

% --- CONFIGURACIÓN DE ENCABEZADO Y PIE DE PÁGINA ---
\pagestyle{fancy}
\fancyhf{} 
\setlength{\headwidth}{\textwidth} % Hace que el ancho del encabezado coincida con el ancho del texto
\setlength{\headheight}{15pt}  % Ajusta la altura del encabezado
\setlength{\headsep}{20pt}     % Ajusta la separación entre el encabezado y el contenido
% Encabezado

\fancyhead[R]{
    \textbf{Curso: 3R4}  \textbf{ - Año: 2026} \\
    \textbf{Análisis de Sistemas y Señales} 
}

% Pie de página
\fancyfoot[L]{\small UTN-FRC}
\fancyfoot[R]{\thepage} % Número de página a la derecha

% Líneas divisorias
\renewcommand{\headrulewidth}{0.4pt}
\renewcommand{\footrulewidth}{0.4pt}



\begin{document}

% --- CARÁTULA ---
\begin{titlepage}
    \centering
    {\scshape\LARGE Universidad Tecnológica Nacional \par}
    \vspace{0.3cm}
    {\scshape\Large Facultad Regional Córdoba \par}
    \vspace{1.5cm}

    {\scshape\Large Análisis de Señales y Sistemas \par}
    \vspace{0.3cm}

    {\Huge\bfseries Trabajo Práctico N°1 \par}
    \vspace{0.5cm}

    {\LARGE  \par}

    \vspace{1cm}

    \begin{flushleft}
    \textbf{Integrantes:} \\
    \quad Rodas Angel -- Legajo 95024 \\
    \quad Dominguez Matías -- Legajo 400462 \\
    \quad Rivera Pedro -- Legajo 400716\\

    \vspace{0.5cm}

    \textbf{Curso:} 2R4 \\
    \textbf{Docentes:} \\
    \quad Ing. Mario R. Modesti  \\
    \quad Ing. Dimas  A. Benasulin \\
    \textbf{Año:} 2026 \\
    \end{flushleft}

    \vfill
\end{titlepage}

% --- ÍNDICE ---
\begin{spacing}{1.5} 
    % Si NO querés encabezado en el índice, descomentá la siguiente línea:
     \thispagestyle{plain} 
    \tableofcontents
\end{spacing}
\newpage

\section{Ejercicio 1}

Considerar la siguiente función en variable compleja $f(z) = \sen \left(e^{j\pi/2}z^2\right)$

\vspace{0.4cm}

$a)$ Obtener el conjunto de ceros de la función $z \in S : f(z) = 0$.

\vspace{0.4cm}

$b)$ Demostrar que $S$ es un conjunto, más aún, hay una cantidad infinitas de ceros de la función sobre las siguientes rectas parametrizadas:
\[
S \cap \left\{t\, e^{-j\pi/4} : t \in \mathbb{R}\right\} = S_1
\]
\vspace{0.2cm}
\[
S \cap \left\{t j\, e^{-j\pi/4} : t \in \mathbb{R}\right\} = S_2
\]

\vspace{0.4cm}

$c)$ El conjunto $S_1$ es discreto, hacer un gráfico de módulos y otro gráfico de argumentos
principal.

\vspace{0.4cm}

$d)$ El conjunto $S_2$ es discreto, hacer un gráfico de módulos y otro gráfico de argumentos
principal.


\subsection{a) Obtener el conjunto de ceros de la función $z \in S : f(z) = 0$. }
Para arrancar, recordemos que la expresión $e^{j\pi/2}$  es igual a $j$:
\[ e^{j\pi/2} = \cos(\pi/2) + j\sin(\pi/2) = j \]
Ahora, reescribiendo la función:
\[ f(z) = \sin(e^{j\pi/2} z^2) \]
\[ f(z) = \sin(j z^2) \]



Igualamos la función a cero para buscar las raíces:
\[ \sin(j z^2) = 0 \]
Sabemos que el seno se hace cero cuando su argumento es un múltiplo entero de $\pi$:
\[ j z^2 = k\pi \quad \text{para } k \in \mathbb{Z} \]
Si despejamos $z^2$ pasando la $j$ dividiendo, nos queda:
\[ z^2 = \frac{k\pi}{j} = -j k\pi \]
\newpage

Para obtener $z$, tenemos que analizar qué pasa con las raíces dependiendo del signo que tenga 
el entero $k$:
\begin{itemize}
    \item $k > 0$: Acá el término $-jk\pi$ es un número imaginario puro negativo. 

    \vspace{0.1cm}
    
    Pasándolo a su forma polar se escribe como $k\pi e^{-j\pi/2}$. Aplicando la raíz cuadrada:
    \[ z = \sqrt{k\pi e^{-j\pi/2}} = \pm \sqrt{k\pi} e^{-j\pi/4} \]
    
    \item $k < 0$: En este caso el término quedaría $-j(-k)\pi = jk\pi$. 

    \vspace{0.1cm}
    
    Siendo este último un imaginario puro positivo. En forma polar sería $k\pi e^{j\pi/2}$.
    \[ z = \sqrt{k\pi e^{j\pi/2}} = \pm \sqrt{k\pi} e^{j\pi/4} \]
    
    \item $k = 0$: Esto es directo.
    \[ z^2 = 0 \implies z = 0 \]
\end{itemize}

Juntando todo esto, el conjunto total de ceros $S$ queda definido:
\[ S = \left\{ \pm \sqrt{k\pi} e^{-j\pi/4} : k \in \mathbb{N} \right\} \cup \left\{ \pm \sqrt{n\pi} e^{j\pi/4} : n \in \mathbb{N} \right\} \cup \{0\} \]

\vspace{0.5cm}

\subsection{b) Demostrar que hay infinitos ceros sobre las rectas.}

Para $S_1$: \qquad $z(t) = t e^{-j\pi/4}$ con la restricción $t \in \mathbb{R}$. 

\vspace{0.4cm}

Reemplazamos esta parametrización  en nuestra función $f(z) = \sin(jz^2)$:
\begin{align*}
    f(z(t)) &= \sin \left( j (t e^{-j\pi/4})^2 \right) \\
            &= \sin \left( j \cdot t^2 e^{-j\pi/2} \right)
\end{align*}	
	
Sabemos que $e^{-j\pi/2} = -j$, la expresión queda:
\[ f(z(t)) = \sin \left( j \cdot t^2 (-j) \right) = \sin(-j^2 t^2) = \sin(t^2) \]
Para que $z(t)$ sea efectivamente un cero de la función, exigimos que dé cero:
\[ \sin(t^2) = 0 \implies t^2 = k\pi \quad \text{con } k \in \mathbb{Z} \]
Como $t$ es un número real, su cuadrado no puede ser negativo, así que $k$ tiene que ser mayor o igual a cero. Despejando $t$:
\[ t = \pm \sqrt{k\pi} \quad \text{para } k \ge 0, 1, 2, 3, \dots \]
Como tenemos infinitos valores enteros posibles para $k \ge 0$, también vamos a tener infinitos valores reales para $t$. Con esto comprobamos que \textbf{hay infinitos ceros sobre la recta $S_1$}.

\vspace{0.4cm}

Para $S_2$: \qquad $z(t) = tje^{-j\pi/4}$ con $t \in \mathbb{R}$.

\vspace{0.4cm}

Sabemos que $j = e^{j\pi/2}$, si reescribimos la parametrización queda:

\[
z(t) = te^{j\pi/2} e^{-j\pi/4} = te^{j\pi/4}
\]

Reemplazamos en la función:

\begin{align*}
f(z(t)) &= \sin\left(j(te^{j\pi/4})^2\right) \\
&= \sin\left(j \cdot t^2 e^{j\pi/2}\right)
\end{align*}

Teniendo en cuenta nuevamente que $e^{j\pi/2} = j$, nos queda:

\[
f(z(t)) = \sin\left(j \cdot t^2(j)\right) = \sin(j^2t^2) = \sin(-t^2) = -\sin(t^2)
\]

Igualamos a cero para buscar las raíces:

\[
-\sin(t^2) = 0 \implies \sin(t^2) = 0 \implies t^2 = k\pi \quad \text{con } k \in \mathbb{N}_0
\]

Y llegamos otra vez a la misma condición para el parámetro:

\[
t = \pm\sqrt{k\pi} \quad \text{para } k = 0, 1, 2, 3, \dots
\]

Igual que en el caso anterior, como $k$ puede tomar infinitos valores naturales, demostramos que hay \textbf{infinitas soluciones} que caen sobre esta segunda recta $S_2$.



\subsection{c) Gráficos de módulos y argumentos para $S_1$}

 $S_1$ está formado por los puntos $z_k = \pm \sqrt{k\pi} e^{-j\pi/4}$. Para graficar esto como un espectro discreto, usamos un índice entero $n \in \mathbb{Z}$ que adopte valores positivos como negativos. Lo definimos por partes de la siguiente manera:
\[
z_n = \begin{cases} 
\sqrt{n\pi} \, e^{-j\pi/4} & \text{para } n > 0 \\
0 & \text{para } n = 0 \\
-\sqrt{|n|\pi} \, e^{-j\pi/4} & \text{para } n < 0 
\end{cases}
\]

\vspace{0.3cm}
\noindent Gráfico de Módulos ($|z_n|$ vs $n$): \\
El módulo es la distancia al origen, $|z_n| = \sqrt{|n|\pi}$. En el gráfico, los puntos forman una curva simétrica que crece hacia ambos lados.

\begin{center}
\begin{tikzpicture}
\begin{axis}[
    width=11cm, height=6cm,
    xlabel={Índice $n$},
    ylabel={Módulo $|z_n|$},
    xmin=-5.5, xmax=5.5,
    ymin=0, ymax=4.5,
    xtick={-5,-4,-3,-2,-1,0,1,2,3,4,5},
    axis lines=center,
    grid=both,
    grid style={dashed, gray!30},
    title={Espectro de Módulo para $S_1$ (y $S_2$)}
]
\addplot+ [ycomb, mark=*, thick, blue, samples at={-5,...,5}] {sqrt(abs(x)*pi)};
\end{axis}
\end{tikzpicture}
\end{center}

\vspace{0.3cm}
\noindent Gráfico de Argumentos Principales ($\arg(z_n)$ vs $n$): \\
Tenemos que mirarla definición por partes que armamos:
\begin{itemize}
    \item $n > 0$: El argumento es $-\pi/4$.
    \item $n < 0$: El signo negativo le suma un $\pi$, por lo que el argumento principal queda en $3\pi/4$.
\end{itemize}

\begin{center}
\begin{tikzpicture}
\begin{axis}[
    width=11cm, height=6cm,
    xlabel={Índice $n$},
    ylabel={Argumento $\arg(z_n)$},
    xmin=-5.5, xmax=5.5,
    ymin=-3.5, ymax=3.5,
    ytick={-pi/4, 3*pi/4},
    yticklabels={$-\frac{\pi}{4}$, $\frac{3\pi}{4}$},
    xtick={-5,-4,-3,-2,-1,0,1,2,3,4,5},
    axis lines=center,
    grid=both,
    grid style={dashed, gray!30},
    title={Espectro de Argumento Principal para $S_1$}
]
\addplot+ [ycomb, mark=*, thick, red, samples at={1,...,5}] {-pi/4};
\addplot+ [ycomb, mark=*, thick, red, samples at={-5,...,-1}] {3*pi/4};
\end{axis}
\end{tikzpicture}
\end{center}


\subsection{d) Gráficos de módulos y argumentos para $S_2$}

Para la segunda recta, $S_2$ está formado por los puntos $z_{1_k} = \pm \sqrt{k\pi} e^{j\pi/4}$. Aplicando el índice $n$, definimos:
\[
z_{1_n}= \begin{cases} 
\sqrt{n\pi} \, e^{j\pi/4} & \text{para } n > 0 \\

\qquad0 & \text{para } n = 0 \\
-\sqrt{|n|\pi} \, e^{j\pi/4} & \text{para } n < 0 
\end{cases}
\]

\vspace{0.3cm}
\noindent Gráfico de Módulos ($|z_{1_n}|$ vs $n$): \\
El módulo sigue siendo $|z_{1_n}| = \sqrt{|n|\pi}$. El gráfico es igual al del conjunto anterior.

\vspace{0.3cm}

\noindent Gráfico de Argumentos Principales ($\arg(z_{1_n})$ vs $n$): 
\begin{itemize}
    \item $n > 0$: El argumento es $\pi/4$.
    \item $n < 0$: Al desfasar por el signo negativo y ajustar para que quede dentro del rango principal $(-\pi, \pi]$, el argumento nos queda en $-3\pi/4$.
\end{itemize}

\begin{center}
\begin{tikzpicture}
\begin{axis}[
    width=11cm, height=6cm,
    xlabel={Índice $n$},
    ylabel={Argumento $\arg(w_n)$},
    xmin=-5.5, xmax=5.5,
    ymin=-3.5, ymax=3.5,
    ytick={pi/4, -3*pi/4},
    yticklabels={$\frac{\pi}{4}$, $-\frac{3\pi}{4}$},
    xtick={-5,-4,-3,-2,-1,0,1,2,3,4,5},
    axis lines=center,
    grid=both,
    grid style={dashed, gray!30},
    title={Espectro de Argumento Principal para $S_2$}
]
\addplot+ [ycomb, mark=*, thick, orange, samples at={1,...,5}] {pi/4};
\addplot+ [ycomb, mark=*, thick, orange, samples at={-5,...,-1}] {-3*pi/4};
\end{axis}
\end{tikzpicture}
\end{center}
\newpage
		\section{Ejercicio 2}
Considerar la relación entre las variables complejas $z, w \in \mathbb{C}$:
\[
z = \frac{4w + 8}{w - 1}
\]

\vspace{0.4cm}

$a)$ Despejar la variable $w$ para obtener una función $w = f(z)$, y luego, calcular $\lim_{z \to \infty} f(z)$.

\vspace{0.4cm}

$b)$ Escribir $f(z)$ en su forma rectangular, $f(z) = u(x,y) + j\,v(x,y)$

\vspace{0.4cm}

$c)$ Mapear la región $|z - 4| \leq 1$, por medio de $w = f(z)$ y graficar ambos dominios.

\vspace{0.4cm}

$d)$ Sea $C$ el contorno de la región del inciso anterior, orientado positivamente. Evaluar la integral $\oint_C f(z)\, dz$.


        
	
	\subsection{a) Despeje de w y calculo del limite}
    Partimos de:
	\begin{equation}
		z=\frac{4w+8}{w-1}
	\end{equation}
	Multiplicamos ambos lados por $(w-1)$:
	\[ z(w-1) = 4w+8 \]
	Desarrollando:
	\[ zw - z = 4w +8 \]
	Agrupamos los términos con $w$:
	\[ zw - 4w = 8 + z \]
	Factorizamos:
	\[ w(z-4) = 8+z \]
	Entonces:
	\[ w = \frac{8+z}{z-4} \]
\begin{equation}
	\boxed{f(z) = \frac{8+z}{z-4}, \quad z \neq 4}
\end{equation}
		
	\textbf{Cálculo del límite cuando $z \to \infty$:}
	\[ \lim_{z \to \infty} f(z) = \lim_{z \to \infty} \frac{8+z}{z-4} \]
	Dividiendo numerador y denominador por $z$:
	\[ \lim_{z \to \infty} \frac{\frac{8}{z}+1}{1- \frac{4}{z}} = \frac{0+1}{1-0} = 1 \]
	Por lo tanto:
	\[ \boxed{\lim_{z \to \infty} f(z) = 1 }\]
	
\subsection{b) Forma rectangular de $f(z)$}

Sea $z = x + jy$. Entonces:
\[ f(z) = \frac{8+z}{z-4} = \frac{x+8 + jy}{(x-4) + jy} \]
Multiplicamos numerador y denominador por el conjugado del denominador:
\[ f(z) = \frac{(x+8 + jy) \cdot [(x-4) - jy]}{[(x-4) + jy] \cdot [(x-4) - jy]} \]
Desarrollamos el producto:
\[ (x+8)(x-4) - (x+8)jy + jy(x-4) - j^2y^2 \]
Como $j^2 = -1$, resulta:
\[ (x+8)(x-4)+ y^2 +j[y(x-4)- y(x+8)]  \]
Simplificamos:
\[ (x+8)(x-4) = x^2 +4x- 32 \]
\[ y(x-4) - y(x+8) = y(x-4-x-8) = -12y \]
Por lo tanto, la forma rectangular es $f(z) = u(x,y) + j v(x,y)$ con:
\[\boxed{ u(x,y) = \frac{x^2+y^2 +4x- 32}{(x-4)^2 + y^2} \quad ; \quad v(x,y) = \frac{-12y}{(x-4)^2 + y^2} } \]	
\newpage
\subsection{c) Mapeo de la región $|z-4| \leq 1$}
Partimos de la relación:
\[ z - 4 = \frac{4w+8}{w-1} - 4 = \frac{12}{w-1} \]
Tomando módulo y sustituyendo en $|z-4| \leq 1$:
\[\boxed{ \frac{12}{|w-1|} \leq 1 \implies |w-1| \geq 12 }\]
La imagen es el exterior del círculo de centro $(1, 0)$ y radio $12$ en el plano $W$.

\begin{center}
\begin{tikzpicture}[scale=1.5]
    % Ejes
    \draw[->] (-1,0) -- (7,0) node[right] {$Re(z)$};
    \draw[->] (0,-3) -- (0,3) node[above] {$Im(z)$};
    
    % Marcas
    \foreach \x in {2,4,6} \draw (\x,2pt) -- (\x,-2pt) node[below] {\small $\x$};
    
    % Región sombreada (Interior)
    \fill[pattern=north east lines, pattern color=gray!40] (4,0) circle (1);
    \draw[thick] (4,0) circle (1);
    \fill (4,0) circle (1.5pt) node[below] {\small $4$};
    
    % Etiqueta con flecha
    \draw[<-] (4.7, 0.7) -- (5.5, 1.8) node[right] {\small $|z-4|=1$};
    
    % Texto inferior Plano Z
    \node[below] at (3,-3.5) {La región es el interior: $|z-4| \leq 1$};
    \node[below] at (3,-4.5) {\textbf{Figura 1:} Región en el plano $z$: interior del círculo de centro $4$ y radio $1$.};
\end{tikzpicture}
\end{center}
\vspace{3cm} % Espacio físico real entre los dos dibujos

% --- GRÁFICO 2: PLANO W (ABAJO) ---
\begin{tikzpicture}[scale=0.4]
    % Ejes
    \draw[->] (-14,0) -- (16,0) node[right] {$Re(w)$};
    \draw[->] (0,-14) -- (0,14) node[above] {$Im(w)$};
    
    % Marcas
    \foreach \x in {-10,-5,1,5,10} \draw (\x,4pt) -- (\x,-4pt) node[below] {\small $\x$};
    
    % Sombreado exterior (usando recorte para no ensuciar la página)
    \begin{scope}
        \clip (-13.5,-13.5) rectangle (15.5,13.5);
        \fill[pattern=north east lines, pattern color=gray!40, even odd rule] 
            (-20,-20) rectangle (20,20) (1,0) circle (12);
    \end{scope}
    
 \draw[thick] (1,0) circle (12); 
    \fill (1,0) circle (4pt) node[below left] {$1$};
    % Ecuación con flecha
    \draw[<-] (9.4, 8.4) -- (14, 11) node[right] {\small $|w-1| = 12$ };
    
    % Texto inferior Plano W
    \node[below] at (1,-15) {La imagen es el exterior de la circunferencia};
    \node[below] at (1,-16.5) {\textbf{Figura 2:} Imagen en el plano $w$: exterior del círculo 
    de centro $1$ y radio $12$.};
\end{tikzpicture}

\subsection{d) Cálculo de la integral de contorno}

Partimos de la función obtenida anteriormente:
\begin{equation}
    f(z) = \frac{z+8}{z-4}
\end{equation}

Se desea calcular la integral sobre el contorno $C: |z-4|=1$:
\[ \oint_{C} f(z) \, dz = \oint_{C} \frac{z+8}{z-4} \, dz \]

\textbf{Análisis de la singularidad:}
La función presenta una singularidad en $z=4$. Dado que el contorno $C$ es una circunferencia de radio $1$ centrada en $4$, la singularidad se encuentra dentro del camino de integración.

\textbf{Aplicación del Teorema Integral de Cauchy:} \\
Utilizamos la fórmula integral de Cauchy para una función de la forma $\frac{g(z)}{z-z_0}$:
\[ \oint_{C} \frac{g(z)}{z-z_0} \, dz = 2\pi j \cdot g(z_0) \]

Identificamos los componentes en nuestra integral:
\begin{itemize}
    \item $g(z) = z+8$
    \item $z_0 = 4$
\end{itemize}

Evaluamos $g(z)$ en el punto de la singularidad:
\[ g(4) = 4 + 8 = 12 \]

Sustituyendo en la fórmula:
\[ \oint_{C} \frac{z+8}{z-4} \, dz = 2\pi j \cdot (12) \]

Por lo tanto, el resultado final es:
\begin{equation}
\boxed{    \oint_{C} \frac{z+8}{z-4} \, dz = 24\pi j}
\end{equation}

\newpage

\section{Ejercicio 3}
a) Solucionar la ecuación cuadrática $3z^2 + (2-3j)z - 1 - 3j = 0$, exponiendo claramente su 
metodología.

\vspace{0.4cm}

b) Factorizar el polinomio cuadrático $4z^2 + 12z + 34$, sabiendo que una de sus raices es $z_1
= -\frac{3}{2} + \frac{5}{2}j$, exponiendo claramente su metodología.


\subsection{a) Solucionar la ecuación cuadrática}


Se pide encontrar las soluciones de la siguiente ecuación cuadrática. 
\[3z^2 + (2-3j) z -1-3j= 0 \]    

Para ello, se utiliza la fórmula de Bhaskara.
\[ x_1x_2 = \frac{-b \pm \sqrt{b^2 - 4ac}}{2a}\]
De la cuál se determinan los coeficientes de la ecuación. En este caso, los coeficientes son los siguientes: 
\begin{equation*}
\centering
a= 3 \hspace{2cm}  b= 2-3j \hspace{2cm}  c= -1-3j    
\end{equation*}
\\
Resulta en la ecuación:
\[z_1z_2= \frac{-(2-3j) \pm \sqrt{(2-3j)^2-4(3)(-1-3j)}}{2(3)}\]
Se acomodan los términos y resuelven los paréntesis
\[z_1z_2= \frac{-2+3j}{6} \pm \frac{\sqrt{(4-12j-9)- (-12-36j)}}{6}\]
\[z_1z_2= -\frac{1}{3}+\frac{1}{2}j \pm \frac{\sqrt{7+24j}}{6}\]

Ya casi resuelta la ecuación, debemos encontrar la raíz del número
complejo $7 + 24j$, primero convirtiéndolo en su forma polar y luego 
aplicando la \textbf{Fórmula de De Moivre}.

\[r= \sqrt{7^2+24^2} = \sqrt{625} \Rightarrow \qquad r= 25 \] 
\[\theta = \arctan\frac{24}{7} \Rightarrow \qquad\theta= 0,41\pi\]

\begin{equation*}
    7+24i=25(\cos{0,41\pi} +j\sin{0,41\pi)} \qquad \textbf{Forma Polar}
\end{equation*}
\[\]

\begin{equation*}
    \sqrt[n]{z}=r^n(\cos{\frac{\theta-2k\pi}{n}}+
    j\sin{\frac{\theta-2k\pi}{n}}) \qquad\textbf{Fórmula de De Moivre}
\end{equation*}
\newpage

La aplicación de la fórmula debe considerar las diferentes soluciones 
a la misma raíz, para ello se toman diferentes valores de k, en un 
intervalo de $0 \leq k \leq n-1$ en nuestro caso los valores son 
$k= 0$ y $k=1$. Sin embargo, luego demostraremos que da igual qué 
valor se tome porque el resultado es el mismo.

\begin{enumerate}
    \item[\textbf{k= 0}] 
    \[\sqrt{7+24j}=\sqrt{25}\cdot\left(\cos{\frac{0,41\pi+2(0)\pi}
            {2}}+j\sin{\frac{0,41\pi+2(0)\pi}{2}}\right)\]
            \[\sqrt{7+24j}= 5 \cdot(0,8+j0,6) \]
            \[\sqrt{7+24j}= 4 + 3j\]

            
    \item[\textbf{k= 1}]
    \[\sqrt{7+24j}=\sqrt{25}\cdot\left(\cos{\frac{0,41\pi+2(1)\pi}
            {2}}+j\sin{\frac{0,41\pi+2(1)\pi}{2}}\right)\]
            \[\sqrt{7+24j}= 5 \cdot(-0,8-j0,6) \]
            \[\sqrt{7+24j}= -4 - 3j\]
\end{enumerate}
En este caso, utilizaremos el valor de $k=0$ para la resolución.
\[z_1z_2= -\frac{1}{3}+\frac{1}{2}j \pm (\frac{4}{6}+\frac{3}{6}j)\]
\[z_1z_2= -\frac{1}{3}+\frac{1}{2}j \pm (\frac{2}{3}+\frac{1}{2}j)\]



\begin{equation}
        z_1= -\frac{1}{3}+\frac{1}{2}j + \frac{2}{3}+\frac{1}{2}j
        \label{ecuacionz1}
\end{equation}

\begin{equation*}
    \boxed{z_1=\frac{1}{3}+j}
\end{equation*}

\begin{equation}
        z_2= -\frac{1}{3}+\frac{1}{2}j - (\frac{2}{3}+\frac{1}{2}j)
        \label{ecuacionz2}
\end{equation}

\begin{equation*}
    \boxed{z_2=-1}
\end{equation*} 


La solución final a la ecuación $3z^2 + (2-3j) z -1-3j= 0$ es la siguiente:
\begin{equation*}
    \boxed{3\cdot[z-(\frac{1}{3}+j)]\cdot[z-(-1)]=0}
\end{equation*}

En el caso de haber utilizado la raíz de $k=1$ lo que hubiera sucedido es que, debido a que
es el número complejo opuesto al de la raíz de $k=0$ este se ve contrarrestado por el efecto
del $\pm$ en la ecuación. Esto es lo que sucede:
\[\pm \frac{(-4-3j)}{6} \Rightarrow + \frac{(-4-3j)}{6}=-\frac{2}{3}-\frac{1}{3}j\]
\[\pm \frac{(-4-3j)}{6}\Rightarrow-\frac{(-4-3j)}{6}=\frac{2}{3}+\frac{1}{3}j\]
Se observa que el resultado del término en sus distintos signos es igual al obtenido
en la ecuación \eqref{ecuacionz1} y la ecuación \eqref{ecuacionz2}.

\subsection{b) Factorizar el polinomio cuadrático}
\textbf{Identificación de raíces:} \\
Dado que el polinomio $P(z) = 4z^2 + 12z + 34$ tiene coeficientes puramente
reales ($a=4$, $b=12$ y $c=34$), sabemos
que si un número complejo $z_1$ es raíz de la ecuación polinomial, entonces
su complejo conjugado $z_2 = \overline{z_1}$ también es otra raíz.

El enunciado nos proporciona la primera raíz:
\[
z_1 = -\frac{3}{2} + \frac{5}{2}j
\]
Por lo tanto, deducimos inmediatamente que la segunda raíz es su conjugado:
\[
z_2 = \overline{z_1} = -\frac{3}{2} - \frac{5}{2}j
\]

\textbf{Construcción de la forma factorizada:} \\
Todo polinomio cuadrático de la forma $az^2 + bz + c$ con raíces $z_1$ y $z_2$ se puede factorizar expresándolo como:
\[
P(z) = a(z - z_1)(z - z_2)
\]
Sustituyendo nuestro coeficiente principal $a = 4$ y las raíces $z_1$ y $z_2$ obtenemos la ecuación factorizada final.
\begin{equation*}
\boxed{P(z)=4 \left[ z - \left(-\frac{3}{2} + \frac{5}{2}j\right) \right] \left[ z - \left(-\frac{3}
{2} - \frac{5}{2}j\right) \right] }
\end{equation*}

\newpage
\section{Ejercicio 4}
b) Evaluar la integral
\[
\int_{0}^{2\pi} \frac{1}{\cos(\theta) + 2\sen(\theta) + 3} \, d\theta
\]
exponiendo claramente su metodología.

\vspace{0.4cm}

c) Evaluar la integral
\[
\int_{-\infty}^{\infty} \frac{2x^2 - 1}{x^4 + 5x^2 + 4} \, dx
\]
exponiendo claramente su metodología.
\vspace{0.4cm}

d) Evaluar la integral
\[
\int_{0}^{\infty} \frac{x \sen(x)}{x^4 + 4} \, dx
\]
exponiendo claramente su metodología.
\subsection{a) Evaluar la integral trigonométrica
}
\[\int_{0}^{2\pi} \frac{1}{\cos(\theta) + 2\sin(\theta) + 3} \, d\theta\]
\textbf{Transformación al plano complejo:} \\
Sea el contorno $C: |z|=1$. Utilizamos las siguientes sustituciones:
\begin{itemize}
    \item $\cos(\theta) = \frac{z + z^{-1}}{2} = \frac{z^2+1}{2z}$
    \item $\sin(\theta) = \frac{z - z^{-1}}{2j} = \frac{z^2-1}{2jz}$
    \item $d\theta = \frac{dz}{jz}$
\end{itemize}

Sustituyendo en la integral original:
\[
\oint_{C} \frac{1}{\left[ \frac{z^2+1}{2z} \right] + 2\left[ \frac{z^2-1}{2jz} \right] + 3} \frac{dz}{jz}
\]

\textbf{Simplificación del denominador:} \\
Introducimos el término $jz$ del diferencial dentro del corchete del denominador:
\[
\oint_{C} \frac{1}{jz \left[ \frac{z^2+1}{2z} + \frac{z^2-1}{jz} + 3 \right]} dz = \oint_{C} \frac{1}{\frac{j(z^2+1)}{2} + (z^2-1) + 3jz} dz
\]
Multiplicamos numerador y denominador por 2 para eliminar la fracción:
\[
\oint_{C} \frac{2}{j z^2 + j + 2z^2 - 2 + 6jz} dz
\]
Agrupando términos por potencias de $z$, obtenemos la función $f(z)$:
\[\boxed{
\oint_{C} \frac{2}{(2+j)z^2 + 6jz + (-2+j)} dz}
\]

\subsection*{Cálculo de Polos}
Buscamos las raíces del denominador $(2+j)z^2 + 6jz + (-2+j) = 0$ mediante Bhaskara:
\[
z_{1,2} = \frac{-6j \pm \sqrt{(6j)^2 - 4(2+j)(-2+j)}}{2(2+j)}
\]
Operando dentro de la raíz: $(2+j)(-2+j) = -4+2j-2j+j^2 = -4-1 = -5$.
\[
z_{1,2} = \frac{-6j \pm \sqrt{-36 - 4(-5)}}{4+2j} = \frac{-6j \pm \sqrt{-16}}{4+2j} = \frac{-6j \pm 4j}{4+2j}
\]

\textbf{Raíces obtenidas:}
\begin{enumerate}
    \item ${z_1 = \frac{-2j}{4+2j} = \frac{-j}{2+j} \cdot \frac{2-j}{2-j} = \frac{-2j+j^2}{4+1} = -\frac{1}{5} - \frac{2}{5}j}$
    \item ${z_2 = \frac{-10j}{4+2j} = \frac{-5j}{2+j} \cdot \frac{2-j}{2-j} = \frac{-10j+5j^2}{4+1} = -1 - 2j}$


    
\end{enumerate}
\begin{equation*}
    \boxed{|z_1| \approx 0,447 \quad (\text{Interior})} \qquad    \boxed{|z_2| \approx 2,236 \quad (\text{Exterior})}
\end{equation*}

\subsection*{Cálculo del Residuo}
Calculamos el residuo para el polo simple $z_1$:
\[
\text{Res}(f, z_1) = \lim_{z \to z_1} (z-z_1) \frac{2}{(2+j)(z-z_1)(z-z_2)} = \frac{2}{(2+j)(z_1-z_2)}
\]
Sustituimos la diferencia de las raíces:
\[
z_1 - z_2 = \left( -\frac{1}{5} - \frac{2}{5}j \right) - (-1 - 2j) = \frac{4}{5} + \frac{8}{5}j = \frac{4}{5}(1+2j)
\]
Calculamos el denominador del residuo:
\[
(2+j) \cdot \frac{4}{5}(1+2j) = \frac{4}{5}(2 + 4j + j + 2j^2) = \frac{4}{5}(2 + 5j - 2) = \frac{4}{5}(5j) = 4j
\]
Por lo tanto:
\[\boxed{
\text{Res}(f, z_1) = \frac{2}{4j} = \frac{1}{2j}}
\]

\subsection*{Cálculo de la Integral Final}
Aplicando el Teorema del Residuo de Cauchy:
 \[
\oint f(z) \, dz = 2\pi j \cdot \text{Res}(f, z_1) = 2\pi j \cdot \left( \frac{1}{2j} \right) = \pi
\] 
\[
\boxed{\oint f(z) \, dz = \pi}
\]

\subsection{b) Evaluar la integral impropia racional}
\[\int_{-\infty}^{\infty} \frac{2x^2 - 1}{x^4 + 5x^2 + 4} \, dx\]

\textbf{Transformación al plano complejo:} \\
Consideramos la función de variable compleja asociada $f(z)$ y tomamos un contorno cerrado $C$ que consiste en el intervalo $[-R, R]$ sobre el eje real y una semicircunferencia $C_R$ de radio $R$ en el semiplano superior:
\[
f(z) = \frac{2z^2 - 1}{z^4 + 5z^2 + 4}
\]
Sabiendo que al hacer $R \to \infty$ la integral sobre $C_R$ tiende a cero, podemos resolver la integral impropia aplicando el Teorema de los Residuos.

\textbf{Factorización del denominador:} \\
Buscamos expresar el denominador en sus factores primos para hallar las singularidades (polos):
\[
z^4 + 5z^2 + 4 = (z^2 + 1)(z^2 + 4) = (z-i)(z+i)(z-2i)(z+2i)
\]

\subsection*{Cálculo de Polos}
Los polos de la función se encuentran donde el denominador se anula. Igualando a cero, obtenemos los siguientes polos simples:
\begin{enumerate}
    \item $z_1 = i \quad (\text{Interior al contorno, semiplano superior})$
    \item $z_2 = 2i \quad (\text{Interior al contorno, semiplano superior})$
    \item $z_3 = -i \quad (\text{Exterior, semiplano inferior})$
    \item $z_4 = -2i \quad (\text{Exterior, semiplano inferior})$
\end{enumerate}
\begin{center}
\begin{tikzpicture}[
    % Definimos el estilo para poner flechas a lo largo del contorno
    flechas/.style={
        decoration={
            markings,
            mark=at position 0.25 with {\arrow[scale=1.5]{latex}},
            mark=at position 0.75 with {\arrow[scale=1.5]{latex}}
        },
        postaction={decorate}
    }
]

    % Dibujamos los ejes cartesianos (x e y)
    \draw[->] (-4,0) -- (4,0) node[right] {$x$};
    \draw[->] (0,-1) -- (0,3.5) node[above] {$y$};

    % Contorno semicircular (arco)
    \draw[flechas, line width=1.2pt, red!70!black] 
        (3,0) arc (0:180:3) 
        node[pos=0.15, above right, text=black] {$C_R$};

    % Segmento sobre el eje real
    \draw[flechas, line width=1.2pt, red!70!black] 
        (-3,0) -- (3,0);

    % Etiquetas de -R y R en el eje x
    \node[below] at (-3,0) {$-R$};
    \node[below] at (3,0) {$R$};

    % Marcamos los polos en i y 2i
    \filldraw (0,1) circle (2pt) node[right=3pt] {$i$};
    \filldraw (0,2) circle (2pt) node[right=3pt] {$2i$};

\end{tikzpicture}
\end{center}
\subsection*{Cálculo de los Residuos}
Calculamos el residuo para el polo simple $z_1 = i$:
\[
\text{Res}(f, z_1) = \lim_{z \to i} (z-i) \frac{2z^2 - 1}{(z-i)(z+i)(z^2+4)} = \frac{2(i)^2 - 1}{(i+i)((i)^2+4)}
\]
\[
\text{Res}(f, z_1) = \frac{-2 - 1}{(2i)(-1+4)} = \frac{-3}{(2i)(3)} = -\frac{1}{2i} \]
\[
\boxed{\text{Res}(f, z_1) = \frac{i}{2}}
\]

Calculamos el residuo para el polo simple $z_2 = 2i$:
\[
\text{Res}(f, z_2) = \lim_{z \to 2i} (z-2i) \frac{2z^2 - 1}{(z^2+1)(z-2i)(z+2i)} = \frac{2(2i)^2 - 1}{((2i)^2+1)(2i+2i)}
\]
\[
\text{Res}(f, z_2) = \frac{2(-4) - 1}{(-4+1)(4i)} = \frac{-9}{(-3)(4i)} = \frac{3}{4i} 
\]
\[
\boxed{\text{Res}(f, z_2)= -\frac{3i}{4}}
\]
\subsection*{Cálculo de la Integral Final}
Planteamos la integral sobre el contorno cerrado $C$, la cual se divide en la integral sobre el arco $C_R$ y la integral sobre el segmento real $[-R, R]$:
\[
\oint_C f(z) \, dz = \int_{C_R} f(z) \, dz + \int_{-R}^{R} f(x) \, dx = 2\pi i \sum_{k=1}^{n} \text{Res}(f(z), z_k)
\]
donde $z_k$ denota los polos en el semiplano superior. Sustituimos los residuos obtenidos para $z_1 = i$ y $z_2 = 2i$:
\[
2\pi i \sum_{k=1}^{2} \text{Res}(f(z), z_k) = 2\pi i \left( \frac{i}{2} - \frac{3i}{4} \right) = 2\pi i \left( -\frac{i}{4} \right) = - \frac{2\pi i^2}{4} = \frac{\pi}{2}
\]

Como el grado del denominador de $f(z)$ es dos unidades mayor que el del numerador, podemos demostrar que la integral $\int_{C_R} f(z) \, dz \to 0$ cuando $R \to \infty$. Entonces tenemos:
\[
\text{V.P.} \int_{-\infty}^{\infty} f(x) \, dx = \lim_{R \to \infty} \int_{-R}^{R} f(x) \, dx = 2\pi i \sum_{k=1}^{n} \text{Res}(f(z), z_k)
\]
Finalmente, sustituyendo nuestra función y el valor de la sumatoria de residuos, obtenemos el resultado:
\[\boxed{
\text{V.P.} \int_{-\infty}^{\infty} \frac{2x^2 - 1}{x^4 + 5x^2 + 4} \, dx = \frac{\pi}{2}}
\]

\subsection{c) Evaluar la integral }
\[\int_{0}^{\infty} \frac{x \sin(x)}{x^4 + 4} \, dx\]
\textbf{Uso de la simetría:} \\
Primero observamos que los límites de integración no son de $-\infty$ a $\infty$. Como el integrando es una función par, ya que $\frac{(-x)\sin(-x)}{(-x)^4+4} = \frac{x\sin(x)}{x^4+4}$, podemos reescribir la integral utilizando la simetría par:
\[
\int_{0}^{\infty} \frac{x \sin(x)}{x^4 + 4} \, dx = \frac{1}{2} \int_{-\infty}^{\infty} \frac{x \sin(x)}{x^4 + 4} \, dx
\]

\textbf{Transformación al plano complejo:} \\
Consideramos la función de variable compleja asociada y utilizamos la fórmula de Euler $e^{iz} = \cos(z) + i\sin(z)$ para formar la integral de contorno:
\[
\oint_C \frac{z e^{iz}}{z^4 + 4} \, dz
\]
donde $C$ es el contorno cerrado que consiste en el intervalo $[-R, R]$ sobre el eje real y la semicircunferencia $C_R$ de radio $R$ en el semiplano superior.

\subsection*{Cálculo de Polos}
Definimos $f(z) = \frac{z}{z^4+4}$. Los polos se encuentran donde el denominador se anula:
\[
z^4 + 4 = 0 \implies z^4 = -4
\]
Expresando $-4$ en forma polar como $4e^{i\pi}$, las raíces se obtienen mediante $z_k = \sqrt[4]{4} e^{i(\frac{\pi + 2k\pi}{4})}$ para $k=0,1,2,3$:
\[\boxed{
z_0 = 1+i} \quad \boxed{z_1 = -1+i} \quad \boxed{z_2 = -1-i} \quad \boxed{z_3 = 1-i
}\]
De estas raíces, los únicos polos que se encuentran en el semiplano superior (interior de nuestro contorno $C$) son $z_0 = 1+i$ y $z_1 = -1+i$.

\subsection*{Cálculo de los Residuos}
Para aplicar el límite, primero factorizamos el denominador agrupando sus raíces complejas conjugadas: $z^4+4 = (z^2-2z+2)(z^2+2z+2)$. Las raíces en el semiplano superior son $z_0 = 1+i$ y $z_1 = -1+i$.

Calculamos el residuo para el polo simple $z_0 = 1+i$:
\[
\text{Res}(f(z)e^{iz}, 1+i) = \lim_{z \to 1+i} (z-(1+i)) \frac{z e^{iz}}{(z-(1+i))(z-(1-i))(z^2+2z+2)}
\]
\[
\text{Res}(f(z)e^{iz}, 1+i) = \frac{(1+i) e^{i(1+i)}}{((1+i)-(1-i))((1+i)^2+2(1+i)+2)}
\]
\[
\text{Res}(f(z)e^{iz}, 1+i) = \frac{(1+i) e^{-1+i}}{(2i)(2i+2+2i+2)}\]

\vspace{0.3cm}

\[
\text{Res}(f(z)e^{iz}, 1+i) = \frac{(1+i) e^{-1+i}}{(2i)(4+4i)} 
\]

\vspace{0.3cm}

\[
\text{Res}(f(z)e^{iz}, 1+i) = \frac{(1+i) e^{-1+i}}{8i(1+i)} 
\]

\vspace{0.3cm}

\[
\text{Res}(f(z)e^{iz}, 1+i) =\frac{e^{-1+i}}{8i}
\]

\vspace{0.3cm}

Calculamos el residuo para el polo simple $z_1 = -1+i$:
\[
\text{Res}(f(z)e^{iz}, -1+i) = \lim_{z \to -1+i} (z-(-1+i)) \frac{z e^{iz}}{(z^2-2z+2)(z-(-1+i))(z-(-1-i))}
\]

\[
\text{Res}(f(z)e^{iz}, -1+i) = \frac{(-1+i) e^{i(-1+i)}}{((-1+i)^2-2(-1+i)+2)((-1+i)-(-1-i))}
\]

\[
\text{Res}(f(z)e^{iz}, -1+i) = \frac{(-1+i) e^{-1-i}}{(-2i+2-2i+2)(2i)} 
\]

\[\text{Res}(f(z)e^{iz}, -1+i) = \frac{(-1+i) e^{-1-i}}{(4-4i)(2i)}\]
\[
\text{Res}(f(z)e^{iz}, -1+i) = \frac{(-1+i) e^{-1-i}}
{8i+8}
\]
Factorizando el denominador como $8i+8 = -8i(-1+i)$, podemos simplificar:
\[
\text{Res}(f(z)e^{iz}, -1+i) = \frac{(-1+i) e^{-1-i}}{-8i(-1+i)} \]
\[\text{Res}(f(z)e^{iz}, -1+i) =-\frac{e^{-1-i}}{8i}\]

Sumando los residuos obtenidos:
\[
\sum \text{Res} = \frac{e^{-1+i}}{8i} - \frac{e^{-1-i}}{8i} = \frac{e^{-1}}{8i} (e^i - e^{-i})
\]
Sabiendo que $\sin(1) = \frac{e^i - e^{-i}}{2i}$, reemplazamos $(e^i - e^{-i})$ por $2i\sin(1)$:
\[\boxed{
\sum \text{Res} = \frac{e^{-1}}{8i} (2i \sin(1)) = \frac{\sin(1)}{4e}
}\]
\subsection*{Cálculo de la Integral Final}
Aplicando el Teorema de los Residuos y el Lema de Jordan cuando $R \to \infty$:
\[
\text{V.P.} \int_{-\infty}^{\infty} \frac{x e^{ix}}{x^4 + 4} \, dx = 2\pi i \left( \frac{\sin(1)}{4e} \right) = \frac{\pi \sin(1)}{2e} i
\]
Separando en parte real e imaginaria mediante $e^{ix} = \cos(x) + i\sin(x)$:
\[
\int_{-\infty}^{\infty} \frac{x \cos(x)}{x^4 + 4} \, dx + i \int_{-\infty}^{\infty} \frac{x \sin(x)}{x^4 + 4} \, dx = 0 + i \frac{\pi \sin(1)}{2e}
\]
Igualando las partes imaginarias:
\[
\text{V.P.} \int_{-\infty}^{\infty} \frac{x \sin(x)}{x^4 + 4} \, dx = \frac{\pi \sin(1)}{2e}
\]
Finalmente, multiplicamos por $\frac{1}{2}$ debido a la simetría inicial:
\[
\int_{0}^{\infty} \frac{x \sin(x)}{x^4 + 4} \, dx = \frac{1}{2} \left( \frac{\pi \sin(1)}{2e} \right) 
\]
\[
\boxed{\int_{0}^{\infty} \frac{x \sin(x)}{x^4 + 4} \, dx = \frac{\pi \sin(1)}{4e}}
\]
\end{document}
